\section{Rank shifts and homology of submodules}
\label{sec:rank-shifted-homology}

Let $K$ be an arbitrary field. A rank shift enlarges the tensor axes while retaining
the simplicial degree and its face and degeneracy indices. The resulting
modules contain explicit representations of graph chain complexes. Their
tensor products also give representations of reduced smash products.
Homology of a simplicial module always means the homology of its
normalized chain complex.

\subsection{Fixed tails and normalized chains}

Fix $k\geq 1$ and a tuple $\mathbf c=(c_1,\ldots,c_k)$ of nonnegative
integers. For $p\geq 0$, put
\[
 I_p(\mathbf c)=\prod_{a=1}^k[p+c_a],
 \qquad X_p(\mathbf c)=K^{I_p(\mathbf c)}.
\]
The coordinates $0,\ldots,p$ on each axis form its simplicial part; the
remaining $c_a$ coordinates form a fixed tail. For an order-preserving map
$\alpha:[q]\to[p]$, define
\[
 \alpha^{+c_a}:[q+c_a]\longrightarrow[p+c_a],\qquad
 \alpha^{+c_a}(j)=
 \begin{cases}
   \alpha(j),&0\leq j\leq q,\\
   p+t,&j=q+t,\quad 1\leq t\leq c_a.
 \end{cases}
\]
This is the ordinal sum of $\alpha$ with the identity on a $c_a$-element
ordered set. Define
\begin{equation}\label{eq:shifted-pullback}
 (\alpha^*T)(j_1,\ldots,j_k)
 =T\bigl(\alpha^{+c_1}(j_1),\ldots,
          \alpha^{+c_k}(j_k)\bigr).
\end{equation}
Since $(\alpha\beta)^{+c_a}=\alpha^{+c_a}\beta^{+c_a}$,
these maps make $X(\mathbf c)$ a simplicial $K$-module. For $p\geq1$ its
degree-$p$ faces are $d_0,\ldots,d_p$, and for $p\geq0$ its degree-$p$ degeneracies are
$s_0,\ldots,s_p$. Thus $d_i$ deletes coordinate $i$ on every axis and
$s_i$ duplicates coordinate $i$ on every axis. The tail lengths do not
increase the number of face or degeneracy operators.

Replacing $\mathbf c$ by
$\mathbf c+r\mathbf 1$, where $r\geq0$ and
$\mathbf1=(1,\ldots,1)$, is a common rank shift. We write
\[
 X^{(k,r)}=X(r,\ldots,r),\qquad
 X^{(k,r)}_p=K^{[p+r]^k}
\]
for the equal-tail case. These definitions also include unequal axis
lengths without a common shift. The DSTM of the preceding sections is
the case $\min_a c_a=0$. Allowing all tail lengths to be positive
extends that construction by enlarging its degree-zero module.

For $m=(m_1,\ldots,m_k)\in I_p(\mathbf c)$, let $E_m$ denote the tensor
unit at $m$, and put $\operatorname{val}(m)=\{m_1,\ldots,m_k\}$.
We use the Moore normalization
\[
N_pX=\bigcap_{i=1}^p\ker d_i,\qquad
 \partial=d_0:N_pX\longrightarrow N_{p-1}X\quad(p\geq1).
\]
Here $N_0X=X_0$ and the complex is unaugmented.

\begin{proposition}[Normalized basis and ambient homology]
\label{prop:shifted-normalized-basis}
For $p\geq1$, the following sets are bases of the indicated spaces:
\begin{align*}
 N_pX(\mathbf c):&\quad
   \{E_m:\{1,\ldots,p\}\subseteq\operatorname{val}(m)\},\\
 Z_p(NX(\mathbf c)):&\quad
   \{E_m:\{0,\ldots,p\}\subseteq\operatorname{val}(m)\}.
\end{align*}
Moreover,
\begin{equation}\label{eq:shifted-ambient-homology}
 H_p(NX(\mathbf c))=0\quad(p>0),\qquad
 H_0(NX(\mathbf c))
 \cong K^{\prod_{a=1}^k\{1,\ldots,c_a\}}.
\end{equation}
In particular, $\dim_K H_0=\prod_{a=1}^k c_a$, with value zero if
any $c_a=0$.
\end{proposition}

\begin{proof}
The image $d_iE_m$ is zero precisely when some coordinate of $m$ equals
$i$. On the tensor units whose coordinates all avoid $i$, the map $d_i$
is a bijection onto the tensor-unit basis in degree $p-1$: each coordinate
strictly greater than $i$ is decreased by one. Distinct such units
therefore have distinct images. It follows that $\ker d_i$ is spanned by
the units whose index tuples contain $i$. Intersecting these coordinate
subspaces proves both basis statements.

For $p\geq1$, let $W_p$ be the span of the normalized tensor units whose
indices avoid zero. The basis descriptions give
\[
 N_pX=Z_p(NX)\oplus W_p.
\]
On $W_p$, the differential subtracts one from every coordinate. For
$p\geq2$, this is a bijection from its basis onto the basis of
$Z_{p-1}(NX)$. Consequently every positive-degree cycle is a boundary.
For $p=1$, the image $d_0W_1$ is exactly the span of the degree-zero units
whose index tuples contain zero. The complementary degree-zero units
have $1\leq m_a\leq c_a$ for every $a$. Their residue classes give the
asserted basis of $H_0$.
\end{proof}

The vanishing of positive ambient homology does not force the same
vanishing in a simplicial submodule. A cycle in a submodule may have all
of its ambient fillings outside that submodule. The support description
nevertheless imposes a bound depending only on the number of axes.

\begin{corollary}[Bounds for submodule homology]
\label{cor:shifted-submodule-bound}
If $Y\subseteq X(\mathbf c)$ is a simplicial submodule with $k$ tensor
axes, then
\[
 H_p(NY)=0\qquad(p\geq k).
\]
For two axes, $\dim_K H_1(NY)\leq2$, independently of the tail lengths.
\end{corollary}

\begin{proof}
For $p\geq1$,
\[
 Z_p(NY)=Y_p\cap\bigcap_{i=0}^p\ker d_i
         \subseteq Z_p(NX).
\]
A tuple of $k$ coordinates cannot contain the $p+1$ distinct values
$0,\ldots,p$ when $p\geq k$. In degree one with two axes, the only
tuples containing both zero and one are $(0,1)$ and $(1,0)$.
Thus the ambient cycle space has basis $E_{01},E_{10}$, and taking a
subspace and then a quotient cannot increase its dimension.
\end{proof}

\subsection{An embedding of graph chains}

Let $G$ be a finite graph with vertices $1,\ldots,v$, where $v\geq1$,
and edges $1,\ldots,e$. Choose an orientation of each edge, and write
$u_j\to w_j$ for edge $j$. Loops and parallel edges are allowed. Regard
$G$ as the simplicial set with these vertices and nondegenerate edges,
with no nondegenerate simplices in higher degrees. Its realization is
the usual one-dimensional CW complex of the graph.

For $r\geq\max\{v,e\}$, define tensors in $X^{(3,r)}$ by
\begin{align}
 a_t&=E_{0,t,t}\in X^{(3,r)}_0
       &&(1\leq t\leq v),\label{eq:graph-vertices}\\
 T_j&=E_{0,u_j+1,u_j+1}+E_{1,w_j+1,w_j+1}
          +E_{0,1,j+1}\in X^{(3,r)}_1
       &&(1\leq j\leq e).\label{eq:graph-edges}
\end{align}
Deleting index one or zero, respectively, gives
\begin{equation}\label{eq:graph-endpoints}
 d_1T_j=a_{u_j},\qquad d_0T_j=a_{w_j}.
\end{equation}

\begin{theorem}[Graph representation]
\label{thm:graph-representation}
The assignments in \eqref{eq:graph-vertices}--\eqref{eq:graph-edges}
extend to an injective simplicial linear map
\[
 F:K[G]\longrightarrow X^{(3,r)}.
\]
Its image $Y_G$ is the simplicial submodule generated by all $a_t$ and
$T_j$. Its normalized chain complex is isomorphic to
\begin{equation}\label{eq:graph-chain-complex}
 0\longrightarrow K^e
 \xrightarrow{\ \partial\ }K^v\longrightarrow0,
 \qquad \partial(\mathbf e_j)=\mathbf e_{w_j}-\mathbf e_{u_j},
\end{equation}
where $\mathbf e_j$ and $\mathbf e_t$ denote the edge and vertex basis vectors.
If $G$ has $b$ connected components, then
\[
 H_0(NY_G)\cong K^b,\qquad
 H_1(NY_G)\cong K^{e-v+b},\qquad
 H_p(NY_G)=0\quad(p\geq2).
\]
\end{theorem}

\begin{proof}
The endpoint identities \eqref{eq:graph-endpoints} define a simplicial
map from $G$ to the underlying simplicial set of $X^{(3,r)}$; extend it
linearly to obtain $F$. Its image is the stated generated submodule.

The normalized degree-zero map sends the vertex basis to the distinct
tensor units $a_t$, so it is injective. In degree one, a basis of
$N_1K[G]$ is formed by
$j-s_0u_j$. Its image is
\begin{equation}\label{eq:graph-normalized-edges}
 b_j=T_j-s_0a_{u_j}
   =E_{1,w_j+1,w_j+1}-E_{1,u_j+1,u_j+1}+E_{0,1,j+1}.
\end{equation}
The coefficient at $(0,1,j+1)$ is one in $b_j$ and zero in every
$b_{j'}$ with $j'\neq j$. These coordinates also have zero coefficient
in the endpoint terms, since those terms have first coordinate one.
Hence the $b_j$ are linearly independent, including for loops and
parallel edges. The normalized map is therefore injective in degrees
zero and one. It is injective in every higher degree because
$N_pK[G]=0$ there.

The natural direct-sum decomposition into degeneracies of normalized
chains, Theorem~\ref{thm:EZ}, now shows that $F$ is injective in every
simplicial degree. Finally,
\[
 d_0b_j=a_{w_j}-a_{u_j},
\]
which identifies the normalized differential with the incidence map in
\eqref{eq:graph-chain-complex}. On each connected component, the image
of this map is the subspace of vertex coefficients having total sum
zero: the edge differences span this subspace along a spanning tree.
The incidence map consequently has rank $v-b$, over every field.
The homology formulas follow.
\end{proof}

\begin{example}[The Petersen graph]\label{ex:petersen-homology}
For the Petersen graph, $v=10$, $e=15$, and $b=1$. The construction with
$r=15$ gives a simplicial submodule of $X^{(3,15)}$ with
\[
 H_0\cong K,\qquad H_1\cong K^6,
 \qquad H_p=0\quad(p\geq2).
\]
Corollary~\ref{cor:shifted-submodule-bound} shows that no simplicial
submodule of a two-axis DSTM can have this first homology, regardless of
its rank shift.
\end{example}

\subsection{Tensor products and reduced smash products}

For tuples $\mathbf c=(c_1,\ldots,c_k)$ and
$\mathbf d=(d_1,\ldots,d_\ell)$, denote their concatenation by
$(\mathbf c,\mathbf d)$. Tensor products of simplicial modules in this
subsection are taken degreewise.

\begin{proposition}[Concatenation of axes]
\label{prop:shifted-tensor-product}
There is a natural isomorphism of simplicial $K$-modules
\begin{equation}\label{eq:axis-concatenation}
 X(\mathbf c)\otimes_K X(\mathbf d)
       \cong X(\mathbf c,\mathbf d),\qquad
 E_m\otimes E_n\longmapsto E_{(m,n)}.
\end{equation}
If $A\subseteq X(\mathbf c)$ and $B\subseteq X(\mathbf d)$ are
simplicial submodules, this isomorphism restricts to an embedding
$A\otimes_KB\hookrightarrow X(\mathbf c,\mathbf d)$.
\end{proposition}

\begin{proof}
In every degree, the tensor-unit bases give the displayed linear
isomorphism. Equation~\eqref{eq:shifted-pullback} shows that it commutes
with every simplicial map: the same order map acts on every axis.
Tensoring inclusions of vector spaces over a field gives an inclusion,
so the restriction is injective in every degree.
\end{proof}

Different tail lengths need not be equalized for
\eqref{eq:axis-concatenation}. If an equal-tail ambient module is
preferred, there are natural zero-padding inclusions
\[
 X(\mathbf c)\hookrightarrow X(\mathbf c')
 \qquad(c_a\leq c'_a\text{ for all }a).
\]
In degree $p$, retain the coordinates $0,\ldots,p+c_a$ on axis $a$ and
set every entry using one of its new tail coordinates to zero. The
fixed-tail maps preserve old and new tails separately, so this padding
commutes with faces and degeneracies. In particular,
\[
 X^{(k,r)}\otimes_KX^{(\ell,r)}\cong X^{(k+\ell,r)}.
\]

For a pointed simplicial set $S$, write
$\widetilde K[S]=K[S]/K[*]$ for its reduced free simplicial module.
The smash product is the degreewise quotient
$S\wedge T=(S\times T)/(S\vee T)$.
An embedding of $K[S]$ also supplies an embedding of
$\widetilde K[S]$: the latter is naturally isomorphic to the kernel of
the augmentation $\varepsilon:K[S]\to K$ by
\[
 [\sigma]\bmod K[*]\longmapsto[\sigma]-[*].
\]
Here $*$ denotes the basepoint simplex in the appropriate degree.
In particular, choosing a vertex as basepoint in
Theorem~\ref{thm:graph-representation} yields such a reduced embedding.

\begin{proposition}[Reduced smash representation]
\label{prop:shifted-smash}
There is a natural isomorphism
\begin{equation}\label{eq:reduced-smash-linearization}
 \widetilde K[S\wedge T]
   \cong\widetilde K[S]\otimes_K\widetilde K[T].
\end{equation}
Consequently, embeddings
\[
 \widetilde K[S]\hookrightarrow X(\mathbf c),\qquad
 \widetilde K[T]\hookrightarrow X(\mathbf d)
\]
induce an embedding
$\widetilde K[S\wedge T]\hookrightarrow X(\mathbf c,\mathbf d)$.
The Eilenberg--Zilber maps give a natural chain homotopy equivalence
\begin{equation}\label{eq:shifted-eilenberg-zilber}
 N\widetilde K[S\wedge T]
   \simeq
 N\widetilde K[S]\otimes_K N\widetilde K[T],
\end{equation}
where the right side is the tensor product of chain complexes.
\end{proposition}

\begin{proof}
In each degree, a basis of either side of
\eqref{eq:reduced-smash-linearization} is indexed by pairs of
nonbasepoint simplices. A pair containing a basepoint represents zero
on both sides. This identification commutes with all simplicial maps,
including those that send a nonbasepoint simplex to the basepoint.
The embedding follows from Proposition~\ref{prop:shifted-tensor-product}.
Applying the shuffle and Alexander--Whitney maps, transported between
the normalized quotient and Moore normalization, gives
\eqref{eq:shifted-eilenberg-zilber}; see
\cite[\S8.5]{Weibel}.
\end{proof}

\begin{example}[The smash product of two circles]
\label{ex:shifted-circle-smash}
In the unshifted two-axis module $X^{(2,0)}$, let
$a=E_{01}\in X^{(2,0)}_1$. Both faces of $a$ are zero. The simplicial
submodule $A$ generated by $a$ is isomorphic to
$\widetilde K[\Delta[1]/\partial\Delta[1]]$. Indeed, sending the
nondegenerate edge to $a$ defines a map from this reduced free module;
its normalized map is an isomorphism onto $Ka$ in degree one and has
zero source in every other degree. The degeneracy decomposition proves
injectivity, and the generated-submodule definition proves surjectivity.
Thus
\[
 N_1A=Ka,\qquad N_pA=0\quad(p\neq1).
\]
Proposition~\ref{prop:shifted-smash} embeds $A\otimes_K A$ in
$X^{(4,0)}$. A normalized two-cycle representing its homology generator
is the shuffle
\begin{equation}\label{eq:circle-smash-shuffle}
 z=(s_1a)\otimes(s_0a)-(s_0a)\otimes(s_1a).
\end{equation}
The simplicial identities give $d_0z=d_1z=d_2z=0$. On this cycle, the
$(1,1)$-component of the Alexander--Whitney map is computed by
$d_2\otimes d_0$ and sends $z$ to $a\otimes a$.
This is the nonzero generator of the tensor product complex
$NA\otimes_KNA$, which is concentrated in degree two. The
Eilenberg--Zilber equivalence therefore gives
\[
 H_2\bigl(N(A\otimes_K A)\bigr)\cong K,
 \qquad H_p\bigl(N(A\otimes_K A)\bigr)=0\quad(p\neq2).
\]
These are the reduced homology groups of the simplicial smash product
of two circles. The computation holds in characteristic two as well,
where the minus sign in \eqref{eq:circle-smash-shuffle} becomes a plus
sign.
\end{example}

\subsection{Sparse computations}

The accompanying Python implementation stores rational coefficients at
nonzero tensor coordinates, together with the simplicial degree and the
fixed tail lengths. Faces, degeneracies, normalization, and concatenation
of axes are computed directly on these coordinates. Exact rational
elimination on the occurring coordinates gives the following dimensions
and ranks, where $\partial_p=d_0:N_p\to N_{p-1}$:
\begin{center}
\begin{tabular}{lrrrrr}
\hline
Representation & $\dim N_0$ & $\dim N_1$ & $\dim N_2$
 & $\rank\partial_1$ & $\rank\partial_2$\\
\hline
Petersen graph & 10 & 15 & 0 & 9 & 0\\
Reduced circle smash & 0 & 1 & 2 & 0 & 1\\
\hline
\end{tabular}
\end{center}
Both computations therefore agree with the preceding homology
calculations. At shift $r=15$, each degree-one graph tensor belongs to an
ambient space with $17^3=4913$ coordinates but has only three nonzero
entries. This sparsity does not give a general efficiency bound:
duplicating one coordinate on each of $k$ axes can replace a single entry
by $2^k$ entries. Source code and reproducible examples are available in
\url{https://github.com/flengyel/SimplicialTensors}; the experiment is
\texttt{experiments/rank\_shifted\_homology.py}.
